\section{Conclusion and Future Work}
\label{sec:conclusion}

In this work, we presented \textbf{PAC-Bayes Merge}, a novel, mathematically grounded framework that provides the first formal, information-theoretic foundation for trajectory-regularized model merging. Addressing the Overfitting-Optimizer Paradox that plagues parameter-space optimization under extreme calibration data scarcity, our framework projects high-dimensional layer-wise ensembling coefficients onto a continuous, low-degree polynomial trajectory. By modeling the trajectory coefficients as the mean parameters of a randomized isotropic Gaussian posterior and defining a spherical Gaussian prior centered at the uniform ensembling baseline, we proved that minimizing Alquier's linear PAC-Bayesian bound analytically justifies a smooth quadratic $L_2$ Consensus-Pulling penalty centered at the stable uniform consensus.

Through extensive evaluation on genuine physical image classification datasets (MNIST, FashionMNIST, CIFAR-10, SVHN) projected via Johnson-Lindenstrauss embeddings into a 14-layer deep MLP residual backbone across 15 random seeds, we demonstrated that PAC-Bayes Merge successfully circumvents transductive overfitting and functional collapse. Our advanced non-isotropic Fisher-guided PAC-Bayes-FIM Merge achieves a Joint Mean accuracy of \textbf{36.13\%}, significantly outperforming the Static Uniform baseline (33.57\%), properly tuned Ties-Merge (29.68\%), DARE-Merge (33.24\%), and unconstrained layer-wise optimization (36.09\%), illustrating that our smooth $L_2$ regularization effectively preserves continuous representative capacity across intermediate network layers without forcing artificial coordinate sparsity. Furthermore, we established a formal equivalence between uniform weight-space merging and Stochastic Weight Averaging (SWA) under SGD sampling noise, and discussed its limitations in disparate multi-task environments where experts reside in structurally distinct basins.

Finally, because our optimized trajectory parameters are statically compiled into final model weights prior to deployment, PAC-Bayes Merge guarantees zero runtime latency and zero additional memory footprint. This makes it an exceptionally practical solution for real-world resource-constrained applications.

\subsection{Future Directions}
Our theoretical and empirical findings open several promising avenues for future research:
\begin{itemize}
    \item \textbf{Adaptive and Heterogeneous Priors:} Rather than employing isotropic Gaussian priors with fixed spherical covariance, future work could investigate layer-wise adaptive variances based on the empirical Fisher Information Matrix of each network layer. This would allow more sensitive layers to receive tighter regularization while permitting less sensitive layers to adjust more freely.
    \item \textbf{Alternative Prior Configurations:} Exploring non-Gaussian priors (such as mixture models or matrix-variate distributions) could provide alternative learning-theoretic regularizers that are more aligned with the non-convex geometry of deep networks.
    \item \textbf{Scaling to Large Language Models (LLMs):} While we demonstrated the efficacy of PAC-Bayes Merge in vision architectures, applying trajectory-regularized merging to massive autoregressive language models (e.g., LLaMA or GPT-class models) during parameter-efficient fine-tuning (PEFT) is a critical next step.
\end{itemize}
